\section{Evaluation}\label{sec:eval}

We ask four questions. Do the boundaries refuse the escalations of
\S\ref{sec:threat} without changing state? Does the deployed evaluator's cost
follow the proven step count? What does enforcement cost at each boundary?
Does custody signing over channel accounts work on a public network?

\paragraph{Method.} Every number below comes from a recorded run, and each
claim tag names the result file it is read from; this paper ran no new
experiments. The Soroban runs execute every contract from its wasm in
\code{soroban-env-host} 27.0.1 under enforcing authorization: Perch's account,
\code{WebAuthn} verifier, document compiler and interpreter as deployed on
testnet, Nido's factory and grant signer, and a test dApp. Each agent call is
one \code{AddressWithDelegates} entry carrying the holder's signature, so every
account's \code{\_\_check\_auth} runs. Soroban CPU instructions and memory are
the host's metered costs, the same on every machine and every run, and they
are the costs of a network running that host version; no network was measured.
Agent-host times are wall clock, from one recorded run on an Apple M5 Max
under macOS 26.5.2 with the build profile, scheduling class, power source and
load recorded; Table~\ref{tab:host-cost} gives the method. The testnet result
is a recorded run of a program that creates its own accounts and keys.

\subsection{Refusals}

The attack suite gives every attempt to the agent the owner authorized, and
snapshots state before and after it: at the Soroban boundary the balances of
the account, the recipients, the attacker and the dApp instances in both
assets, every grant (holder key, caveats, allowance, spent, revoked), and the
account's revision, document hash and pending upgrade; at the host the file
tree's hashes and symlink targets, a second server's request count, the key
keeper's signature log and every agent's verdict on two probe files. A refusal
passes only if the snapshots are equal. Controls are allowed calls of the same
shape as the attack next to them.

\begin{table}[t]
\caption{Attack suite: attempts, refusals with state unchanged, and allowed
controls.}\label{tab:attacks}
\small
\setlength\tabcolsep{4pt}
\begin{tabular}{@{}llrrr@{}}
\toprule
\textbf{Boundary} & \textbf{Category} & \textbf{Attempts} & \textbf{Refused} & \textbf{Controls} \\
\midrule
Soroban & self-escalation & 16 & 16 & 0 \\
 & nested-call smuggling & 8 & 8 & 1 \\
 & delegation workarounds & 18 & 18 & 6 \\
 & oversubscription & 3 & 3 & 2 \\
Host & self-escalation & 10 & 10 & 0 \\
 & delegation workarounds & 21 & 21 & 3 \\
\bottomrule
\end{tabular}
\end{table}

In Soroban simulation all 45 attempts are refused with state unchanged and
the 9 controls pass\claim{attack-suite-soroban}; at the host all 31 are, and
its 3 controls pass\claim{attack-suite-host} (Table~\ref{tab:attacks}). The
Soroban attempts include the grant signer's admin functions through the
account and as the grant signer itself; the account's \code{apply\_doc},
\code{execute} and \code{schedule\_upgrade}; documents that wrongly scope an
agent rule to the account or to the grant signer; transfers to the attacker
nested two to five calls deep behind an allowed call, or left out of the
signed tree; contract deployment; \code{approve} and \code{burn} of the
budget asset; arguments of another type with the same bytes (a muxed
recipient, a \code{u64} amount, a \code{u32} for the bytes of its XDR, a
string for a symbol); revoked and rotated-out keys; a replayed entry; and
siblings and a grandchild spending past an ancestor's allowance. Refusals come
from three places: the account (no rule covers the call), the grant signer
(\code{NotHolder}, \code{AccountCall}, \code{ContractCreation},
\code{UnmeteredCall}, \code{NotAllowed}, \code{Overdrawn}, \code{Revoked}, a
failed holder signature), and the host (no entry covers the call, a used nonce, a
re-entry). The host attempts are tool calls that issue or edit grants,
sub-agents using a wider grant than their parent's, revoked agents, paths that
name the work directory and reach outside it, URLs that name the allowed host
and reach another (another port or scheme, an address literal, a trailing
dot, userinfo, redirects), and signing requests for another key, network or
payload kind. These are simulation and host results; the attack suite has not
been run through a local Stellar network or through the relay.

The suite found two things. A muxed address (\code{M\ldots}) of a recipient
pinned by \code{argEq} is refused although it names the same account's key,
since values compare exactly. And it found a defect in the host's wiring: the
host signed only the account's own authorization preimage, not the
delegated-signer preimage a grant-signer holder signs, so an agent whose holder
key the keeper held could not authorize through its grant signer. The host
now signs both after the same network check\claim{host-delegated-preimage}. The
defect failed closed, and it sat in the unverified wiring, where
Table~\ref{tab:tcb} puts the remaining risk.

\paragraph{Workloads.} Two scripted agent sessions with sub-agents run in the
simulation, each call one transaction through the account: \emph{payments},
with a per-call cap of 100 and an allowance of 400, and a payroll sub-agent
pinned to one payee; and \emph{defi-router}, whose calls nest the account's
transfers up to four contexts deep. Every call has the expected outcome, and
every refused one leaves state unchanged. The router session shows the
positional limit of \S\ref{sec:grants}: its calls carry amounts at different
argument positions, so it bounds amounts by its allowance alone.

\subsection{The cost-bound cross-check}

Theorem~\ref{thm:total}'s bound is about a counted model. To test whether the
deployed evaluator follows the model's step count, the harness measures the
grant signer's \code{\_\_check\_auth} on five sweeps, each varying one thing:
the number of caveats; the length of a byte-string argument, compared by an
\code{argEq} and not compared; the number of contexts; and the length of a
\code{contractIn} list. Every point is measured twice on the same stored
grant: passing, with every caveat holding, and failing, with the first caveat
failing on the first context, so the evaluator stops at once. Loading the
grant, verifying the signature and mapping contexts to actions cost the same
in both, so the difference is the evaluator's own work. For every point the
harness writes the grant and actions in the evaluator's form, and the Lean
package evaluates its own counted evaluator and bound on them.

The contract, the native evaluator and the Lean model agree on every verdict.
Over the 28 pairs, the evaluator's work is 96 to 164 CPU instructions per Lean
step, and a least-squares fit through the origin of the instruction
difference on the step difference gives 99.4 instructions per step with
$R^2 = 0.986$\claim{deployed-cost}. The fit is through the origin over
differences, so it prices evaluator work only; the mapping of contexts to
actions, which grows with argument size, is outside it: with an argument that
no caveat compares, the whole check grows from 1.072M to 1.255M instructions
between 0 and 16 KiB while the model's steps stay at 68. Priced at that
constant, the bound covers the whole check above the zero-caveat cost on every
point, by a factor of more than 12, and the margin widens as arguments,
contexts and lists grow, since the bound grows with the product of grant and
action sizes and the measured cost with their sum\claim{bound-looseness}. That
the evaluator's work is proportional to the step count on every input remains
unproven\claim{deployed-cost-general}.

\subsection{Cost on chain}

Table~\ref{tab:signer-cost} gives the grant signer's whole
\code{\_\_check\_auth} on one context, from the caveats sweep: instantiating
the contract, loading the grant, checking the signature, mapping and
evaluating. The zero-caveat check costs 972,392 instructions, which the cost
model counts as zero. Under 4 caveats each further context costs about 97,000
instructions\claim{soroban-check-cost}. Whole agent transactions through the
account cost 2.17M to 2.69M instructions with one context and 5.04M with four
nested contexts, against a limit of 400M per
transaction\claim{soroban-tx-cost}.

\begin{table}[t]
\caption{Grant signer \code{\_\_check\_auth} on one context, every caveat
holding (CPU instructions and memory bytes).}\label{tab:signer-cost}
\small
\begin{tabular}{@{}rrrr@{}}
\toprule
\textbf{Caveats} & \textbf{CPU} & \textbf{Memory} & \textbf{Lean steps} \\
\midrule
0 & 972{,}392 & 416{,}582 & 1 \\
1 & 1{,}024{,}459 & 420{,}766 & 60 \\
4 & 1{,}168{,}746 & 434{,}627 & 188 \\
16 & 1{,}735{,}302 & 489{,}255 & 700 \\
64 & 4{,}003{,}738 & 706{,}951 & 2{,}748 \\
\bottomrule
\end{tabular}
\end{table}

\subsection{Cost in the host}

Table~\ref{tab:host-cost} times one check three ways: native
\code{grant-eval} on values built once; the host-side check without
WebAssembly, converting the host's WIT values natively and evaluating; and
the component through Wasmtime. Each case is a three-caveat grant and an
action it allows; the chain case checks a three-context tree. The component
check takes 0.30 to 1.01\,\textmu s, 15 to 47 times native and about 3 times
the host-side check without WebAssembly, so the price of running the deployed
build in a component is about a factor of 3 over checking natively in the
host; compiling and instantiating the component takes 22.0\,ms, once per
host\claim{host-check-time}. A whole tool call costs more than its check: a
denied fetch 0.68\,\textmu s, a read resolved and checked 37 to
47\,\textmu s, a signature about 10\,\textmu s, and an allowed fetch from a
local server 160\,\textmu s\claim{host-tool-call-time}. The load average was
about 5.5 during the run, and macOS cannot pin a thread to a core, so these
are the medians of one run, not a bound.

\begin{table}[t]
\caption{Per-check time (ns): medians with 95\% confidence intervals over 101
interleaved rounds, after a 200\,ms warm-up per operation, timed in batches of
at least 2\,ms; ratios within a round.}\label{tab:host-cost}
\small
\setlength\tabcolsep{4pt}
\begin{tabular}{@{}lrrrr@{}}
\toprule
\textbf{Case} & \textbf{Native} & \textbf{No Wasm} & \textbf{Component} & \textbf{Ratio} \\
\midrule
files & 12.5 & 182.5 & 579.9 & 46.5$\times$ \\
 & {\scriptsize 12.4--12.6} & {\scriptsize 180.7--183.6} & {\scriptsize 579.1--581.2} & \\
network & 7.6 & 90.4 & 298.6 & 39.6$\times$ \\
 & {\scriptsize 7.5--7.6} & {\scriptsize 89.9--90.7} & {\scriptsize 298.0--299.2} & \\
signing & 15.9 & 127.7 & 391.5 & 24.6$\times$ \\
 & {\scriptsize 15.8--16.0} & {\scriptsize 126.8--128.3} & {\scriptsize 390.8--392.3} & \\
chain (3 ctx) & 67.9 & 341.6 & 1{,}006.7 & 14.9$\times$ \\
 & {\scriptsize 67.4--68.2} & {\scriptsize 341.0--342.8} & {\scriptsize 1{,}003.8--1{,}009.2} & \\
\bottomrule
\end{tabular}
\end{table}

\subsection{The relay}

The relay was run end to end on a local Stellar network
(\file{stellar/quickstart}, core and RPC v29, protocol 28) with the Dfns client
signing through the Dfns mock. It confirmed an allowed transfer and an allowed
call whose tree nests an allowed transfer; it never had the custody key sign a
transfer over its cap or a call whose nested transfer the grant denies; it
refused a tampered entry at simulation without using a sequence number; and
it served two requests at once on two channels\claim{relay-local-network}.
Nido's unchanged TypeScript client relayed to confirmation through
it\claim{relay-frontend}. These are functional tests; the relay has no cost,
throughput or attack-suite measurement.

\subsection{Custody signing over channel accounts}

On Stellar testnet (protocol 29), three fresh accounts were given the custody
key as their only weighted signer (master weight 0, thresholds 1). For both
Dfns payload kinds that sign a Stellar transaction, \code{Hash} and
\code{Transaction}, a classic payment, a Soroban \code{InvokeHostFunction}
and a fee bump of a Soroban call, each sourced from a channel and signed only
by the custody key, were accepted; Horizon's records show every signature
verifying under the custody key. The same payment signed by the channel's
zeroed master key or by an unrelated key, a Soroban transfer signed by an
unrelated key, and a fee bump whose outer signature was the fee source's
zeroed master key were all refused with \code{txBAD\_AUTH} and used no
sequence number\claim{testnet-channel}. The custody key was held by the Dfns
mock, so this shows what the network accepts, not what Dfns will
sign (\S\ref{sec:relay}).

\paragraph{Not measured.} The local Stellar network for the grant signer and
the attack suite; the relay's cost, throughput and attack suite; throughput
against the size of the channel-account swarm; sibling spends submitted at
once through the swarm; and the host's signature authorizing a call through
a grant signer on chain. For concurrent sibling spends the protocol gives an
argument, not a measurement: two transactions whose footprints write the same
ledger entry conflict, and the result must be ``as if'' they were applied in
sequence~\cite{cap63}, so the later one in the protocol's order reads the
earlier one's charge to their common ancestors.
